\chapter{Réalisation et tests}
\label{ch:realisation-tests}
\minitoc
\clearpage

\section{Introduction}

Après la conception, cette partie décrit la réalisation concrète de PharmacieConnect. Elle présente
l'environnement de travail, les technologies utilisées, l'organisation du dépôt et les principales
fonctionnalités développées dans le backend, le tableau de bord administrateur et l'application mobile. La
fin du chapitre est consacrée à l'intégration, à la configuration, aux tests et aux difficultés rencontrées.

\section{Environnement matériel et logiciel}

Le développement de PharmacieConnect a mobilisé plusieurs outils, car le projet combine un backend, une
interface web d'administration et une application mobile.

\begin{table}[H]
    \centering
    \caption{Environnement de développement}
    \begin{tabular}{L{5cm}L{8cm}}
        \toprule
        \textbf{Élément} & \textbf{Description} \\
        \midrule
        Machine de développement & Ordinateur portable Intel Core i7-8750H, 12 threads, 16 Go RAM, stockage local 343 Go. \\
        Système d'exploitation & Linux 6.17 x86\_64. \\
        Éditeur de code & Visual Studio Code. \\
        Backend & Python, FastAPI, Uvicorn, SQLAlchemy, Pydantic, Pytest \\
        Administration web & Node.js, React, Vite, React Router, Axios, Tailwind CSS \\
        Application mobile & Expo, React Native, React Navigation, i18next \\
        Base de données & SQLite en développement local, PostgreSQL selon la configuration \\
        Outils de test & Pytest, Jest, build Vite et tests fonctionnels manuels \\
        \bottomrule
    \end{tabular}
\end{table}

\section{Technologies utilisées}

Les choix techniques découlent directement des besoins du projet. Le backend devait fournir des API
claires et sécurisées. L'administration nécessitait une interface web capable de gérer les données depuis
un navigateur. L'application mobile, de son côté, devait rester simple à utiliser tout en intégrant la
recherche, la carte et la consultation des gardes.

\begin{longtable}{L{4cm}L{10cm}}
    \caption{Technologies utilisées dans le projet}\\
    \toprule
    \textbf{Technologie} & \textbf{Utilisation} \\
    \midrule
    \endfirsthead
    \toprule
    \textbf{Technologie} & \textbf{Utilisation} \\
    \midrule
    \endhead
    FastAPI & Développement de l'API REST, endpoints publics et endpoints protégés. \\
    SQLAlchemy & Modélisation des entités et accès à la base de données. \\
    Pydantic & Validation des données entrantes et structuration des réponses. \\
    JWT & Authentification, jetons d'accès et jetons de rafraîchissement. \\
    React & Construction de l'interface d'administration. \\
    Vite & Outil de développement et de build pour le tableau de bord administrateur. \\
    Axios & Communication entre le tableau de bord administrateur et l'API. \\
    Tailwind CSS & Mise en forme de l'interface d'administration. \\
    Expo & Développement et exécution de l'application mobile. \\
    React Native & Création des écrans et composants mobiles. \\
    i18next & Gestion de l'internationalisation. \\
    OpenStreetMap / Nominatim & Référence cartographique et géocodage lorsque l'enrichissement des coordonnées est nécessaire \cite{nominatim}. \\
    Pytest & Tests automatisés du backend \cite{pytest}. \\
    Jest & Tests automatisés côté application mobile \cite{jest}. \\
    \bottomrule
\end{longtable}

\section{Organisation du dépôt}

Le projet est regroupé dans un même espace de travail. Cette structure a facilité l'intégration entre
l'API, le tableau de bord et l'application mobile, notamment lors des tests de communication entre les
clients et le backend.

\begin{table}[H]
    \centering
    \caption{Organisation générale du dépôt}
    \begin{tabular}{L{5.6cm}L{7.4cm}}
        \toprule
        \textbf{Dossier} & \textbf{Contenu} \\
        \midrule
        \path{backend_pharmacie/} & API FastAPI, modèles SQLAlchemy, services métier, migrations, sécurité et tests. \\
        \path{admin_pharmacie/} & Tableau de bord React/Vite destiné à l'administration. \\
        \begin{minipage}[t]{5.4cm}\ttfamily\small ouerkema-pharmacieconnect-\linebreak 4c94773cce7d/\end{minipage} & Application mobile Expo/React Native destinée aux utilisateurs finaux. \\
        \path{docs/} & Documentation technique liée à l'architecture, l'API, le backend, le web admin et le mobile. \\
        \bottomrule
    \end{tabular}
\end{table}

\section{Réalisation backend}

Le backend constitue la partie centrale de PharmacieConnect. Il reçoit les requêtes des deux interfaces,
applique les règles métier et interagit avec la base de données. Il regroupe également les mécanismes
d'authentification, de validation, d'audit et d'import.

\subsection{Endpoints publics}

Les endpoints publics répondent surtout aux besoins de l'application mobile. Ils donnent accès aux
informations de consultation sans passer par l'espace d'administration.

\begin{itemize}
    \item \texttt{GET /health} pour vérifier l'état du backend ;
    \item \texttt{GET /api/pharmacies} pour consulter les pharmacies avec pagination ;
    \item \texttt{GET /api/pharmacies/search} pour rechercher par texte ;
    \item \texttt{GET /api/pharmacies/nearby} pour rechercher les pharmacies proches d'une position GPS ;
    \item \texttt{GET /api/pharmacies/\{pharmacy\_id\}} pour consulter le détail d'une pharmacie ;
    \item \texttt{GET /api/gardes} pour consulter les gardes d'une date donnée ;
    \item \texttt{GET /api/medicines} pour rechercher dans le catalogue de médicaments ;
    \item \texttt{GET /api/medicines/\{code\_pct\}} pour consulter le détail d'un médicament.
\end{itemize}

Pour la recherche par proximité, l'API utilise les coordonnées GPS envoyées par le client. Elle calcule la
distance avec les pharmacies disponibles et retourne les résultats triés par distance.

\subsection{Authentification et sessions}

L'authentification est gérée côté backend afin de garder le contrôle sur les comptes et les autorisations.
Lorsqu'une connexion réussit, le serveur retourne un jeton d'accès et un jeton de rafraîchissement.

\begin{lstlisting}[caption={Structure simplifiée d'une réponse de connexion}]
{
  "access_token": "[jeton d'acces]",
  "refresh_token": "[jeton de rafraichissement]",
  "token_type": "bearer",
  "expires_in": 1800
}
\end{lstlisting}

Les routes protégées vérifient ensuite le jeton reçu et le rôle de l'utilisateur. Cette vérification est
réalisée côté serveur afin qu'un simple contrôle d'interface ne puisse pas être contourné.

\subsection{Services métier}

La logique principale est placée dans des services. Ce choix évite de concentrer trop de code dans les
routes et rend les traitements plus faciles à tester.

\begin{itemize}
    \item \texttt{AuthService} pour l'inscription, la connexion, la vérification email, le rafraîchissement et la déconnexion ;
    \item \texttt{PharmacyService} pour la gestion et l'import des pharmacies ;
    \item \texttt{GardeService} pour la gestion et l'import des plannings de garde ;
    \item \texttt{MedicineService} pour l'import, la recherche et la consultation des médicaments ;
    \item \texttt{AuditService} pour la journalisation des actions importantes ;
    \item \texttt{CacheService} pour l'utilisation du cache lorsque la configuration le permet.
\end{itemize}

\subsection{Imports CSV}

Les imports CSV ont été ajoutés pour éviter la saisie manuelle d'un grand nombre de lignes. Avant
l'enregistrement, le backend vérifie l'extension, la taille, les colonnes obligatoires, le format des
valeurs et les contraintes métier.

Lorsqu'une ligne est incorrecte, la réponse indique les informations nécessaires pour corriger le fichier.
Pour les médicaments, le code PCT sert d'identifiant pour mettre à jour un enregistrement existant au lieu
de créer un doublon.

\subsection{Audit et indicateurs}

Les actions sensibles sont enregistrées dans les journaux d'audit. Cela concerne notamment les
connexions, les échecs de connexion, les inscriptions, les créations administratives et les imports. Les
données statistiques alimentent le tableau de bord et donnent une vision synthétique de l'activité.

\section{Réalisation du tableau de bord administrateur}

Le tableau de bord administrateur rassemble les fonctions de gestion dans une interface web destinée aux
administrateurs et aux utilisateurs autorisés.

\subsection{Pages principales}

Les principales pages développées sont :

\begin{itemize}
    \item le tableau de bord pour les compteurs, indicateurs et activités récentes ;
    \item la page des pharmacies pour consulter, créer, modifier, supprimer et importer les pharmacies ;
    \item la page des gardes pour gérer les plannings de garde ;
    \item la page des médicaments pour importer et consulter le catalogue ;
    \item la carte pour visualiser les pharmacies géographiquement ;
    \item la page de management pour gérer les assistants et les autorisations ;
    \item la page des journaux d'audit pour consulter les actions enregistrées ;
    \item les pages de profil, paramètres, notifications et langues.
\end{itemize}

\subsection{Gestion des sessions côté administration}

Le tableau de bord utilise un client API centralisé. Ce client ajoute le jeton JWT aux requêtes protégées.
Si le jeton d'accès expire, il tente de renouveler la session. En cas d'échec, l'utilisateur est redirigé
vers la page de connexion.

\section{Réalisation application mobile}

L'application mobile correspond à la partie utilisée par le citoyen. Elle regroupe la recherche, la
consultation des gardes, la carte, le catalogue de médicaments et les paramètres.

\subsection{Recherche et détail des pharmacies}

Depuis l'écran d'accueil, l'utilisateur peut rechercher une pharmacie par texte et consulter les résultats.
La fiche de détail affiche les informations disponibles, comme le nom, l'adresse, le téléphone, le
gouvernorat et les coordonnées lorsque celles-ci existent.

\subsection{Carte et géolocalisation}

L'application peut demander l'autorisation d'utiliser la position de l'utilisateur. Cette position est
envoyée au backend afin de récupérer les pharmacies proches. La carte affiche ensuite les résultats sous
forme de marqueurs.

\subsection{Calendrier des gardes}

Le calendrier donne à l'utilisateur le choix d'une date et affiche les pharmacies de garde correspondantes.
Cette fonctionnalité répond au besoin principal du projet : trouver rapidement une pharmacie disponible à
une date précise.

\subsection{Catalogue de médicaments}

L'écran des médicaments prend en charge la recherche par nom commercial, DCI ou code PCT. La fiche de détail
affiche les informations présentes dans la base, comme le prix public, le tarif de référence et la catégorie
de remboursement.

\subsection{Paramètres, langues et thème}

L'application prend en charge le français, l'anglais et l'arabe. Le choix de l'arabe active l'affichage RTL.
Les paramètres regroupent aussi le thème, les notifications, les favoris et certaines préférences
utilisateur.

\section{Captures d'écran}

Les captures réelles seront ajoutées après export depuis l'administration web et l'application mobile. Les
emplacements suivants conservent la structure du rapport sans utiliser de vues synthétiques.

\visualslot{Vue du tableau de bord administrateur}
{Capture réelle à insérer après exécution de l'interface d'administration.}
{fig:screenshot-admin-dashboard}

\visualslot{Vue de gestion des pharmacies}
{Capture réelle à insérer après exécution de l'interface d'administration.}
{fig:screenshot-admin-pharmacies}

\visualslot{Vue mobile d'accueil et de recherche}
{Capture réelle à insérer depuis l'application mobile.}
{fig:screenshot-mobile-home}

\visualslot{Vue mobile du calendrier des gardes}
{Capture réelle à insérer depuis l'application mobile.}
{fig:screenshot-mobile-gardes}

\visualslot{Vue mobile du catalogue de médicaments}
{Capture réelle à insérer depuis l'application mobile.}
{fig:screenshot-mobile-medicines}

\section{Intégration API}

L'intégration API relie les deux interfaces au backend. Le tableau de bord administrateur utilise Axios,
tandis que l'application mobile possède sa propre configuration d'appel API. Les deux clients consomment
les endpoints FastAPI selon leurs besoins.

\begin{table}[H]
    \centering
    \caption{Principales intégrations API}
    \begin{tabular}{L{4cm}L{5cm}L{4cm}}
        \toprule
        \textbf{Client} & \textbf{Endpoints utilisés} & \textbf{Objectif} \\
        \midrule
        Application mobile & Pharmacies, gardes, médicaments, authentification & Consultation et parcours citoyen \\
        Tableau de bord admin & Administration, imports, indicateurs, audit, authentification & Gestion opérationnelle \\
        Backend & Base de données, services métier, sécurité & Centralisation et validation \\
        \bottomrule
    \end{tabular}
\end{table}

\section{Déploiement et configuration}

La configuration du backend repose sur des variables d'environnement. Cette approche sépare le code des
valeurs qui changent selon l'environnement, comme la base de données ou les clés de sécurité.

\begin{table}[H]
    \centering
    \caption{Variables d'environnement principales}
    \begin{tabular}{L{6.2cm}L{6.8cm}}
        \toprule
        \textbf{Variable} & \textbf{Rôle} \\
        \midrule
        \texttt{DATABASE\_URL} & Chaîne de connexion à la base de données. \\
        \texttt{SECRET\_KEY} & Clé utilisée pour signer les jetons JWT. \\
        \texttt{ALGORITHM} & Algorithme de signature des jetons. \\
        \texttt{ACCESS\_TOKEN\_EXPIRE\_MINUTES} & Durée de validité du jeton d'accès. \\
        \texttt{REFRESH\_TOKEN\_EXPIRE\_DAYS} & Durée de validité du jeton de rafraîchissement. \\
        \texttt{FRONTEND\_URL} & Origine autorisée pour l'interface web. \\
        \texttt{TRUSTED\_HOSTS} & Liste des hôtes autorisés. \\
        \texttt{ENABLE\_REDIS\_CACHE} & Activation ou désactivation du cache Redis. \\
        \bottomrule
    \end{tabular}
\end{table}

Pour la version de validation, le projet est exécuté localement de manière contrôlée : le backend FastAPI
est lancé avec Uvicorn, l'interface d'administration avec le serveur de développement Vite et
l'application mobile avec Expo. Dans cette configuration, l'intégration entre les trois parties peut être
vérifiée avant une mise en production. Pour un déploiement final, la même architecture peut être placée
derrière un serveur HTTPS, avec PostgreSQL comme base de données persistante et des variables
d'environnement séparées pour la production. Les commandes de lancement local sont placées en annexe afin
de ne pas alourdir le corps du rapport.

\section{Tests unitaires, API et manuels}

\subsection{Stratégie de test}

La validation combine des tests automatisés côté backend et des vérifications manuelles sur les
interfaces. Les tests backend portent surtout sur les règles sensibles : authentification, audit,
restrictions régionales, indicateurs et imports.

\begin{table}[H]
    \centering
    \caption{Stratégie de validation}
    \begin{tabular}{L{4cm}L{9cm}}
        \toprule
        \textbf{Type de test} & \textbf{Objectif} \\
        \midrule
        Tests unitaires backend & Vérifier une règle métier isolée dans un service. \\
        Tests API & Vérifier les endpoints, les statuts HTTP, les réponses et les droits d'accès. \\
        Tests d'import & Contrôler les fichiers CSV, les colonnes, les lignes invalides et les bilans d'import. \\
        Tests de sécurité & Vérifier les rôles, la vérification email, les jetons et les accès interdits. \\
        Tests manuels web & Valider les pages d'administration, formulaires, tableaux, cartes et imports. \\
        Tests manuels mobile & Valider la recherche, la carte, le calendrier, les langues, les paramètres et les médicaments. \\
        \bottomrule
    \end{tabular}
\end{table}

\subsection{Résultats de validation automatisée}

Une campagne de vérification a été exécutée sur les trois parties du projet le 13 mai 2026. Elle ne
remplace pas une recette complète en production, mais elle confirme que les modules principaux compilent
ou passent leurs tests automatisés dans l'environnement local.

\begin{longtable}{L{3.6cm}L{4.7cm}L{5.7cm}}
    \caption{Résultats de validation exécutés}\\
    \toprule
    \textbf{Partie vérifiée} & \textbf{Commande} & \textbf{Résultat observé} \\
    \midrule
    \endfirsthead
    \toprule
    \textbf{Partie vérifiée} & \textbf{Commande} & \textbf{Résultat observé} \\
    \midrule
    \endhead
    Backend : indicateurs, assistants et audit & \texttt{pytest} sur les fichiers ciblés & 14 tests réussis, 21 avertissements, durée 4,45 s. \\
    Application mobile & \texttt{npm test -- --runInBand} & 7 tests réussis sur la recherche mobile par localisation, durée 2,781 s. \\
    Administration web & \texttt{npm run build} & Build Vite terminé avec succès, génération du dossier \texttt{dist/}. \\
    Import médicaments & \texttt{pytest} sur le fichier d'upload médicaments avec limite de temps & Exécution interrompue après 20 s sur le premier test ; ce point doit être stabilisé pour la validation complète de ce module. \\
    \bottomrule
\end{longtable}

Le build de l'administration a produit les artefacts \texttt{index.html}, CSS et JavaScript. Vite a signalé
un avertissement concernant une taille de bundle JavaScript supérieure à 500 kB après minification. Cet
avertissement n'empêche pas le build, mais il indique qu'un découpage dynamique des pages pourrait être
ajouté dans une future optimisation.

\subsection{Tests automatisés du backend}

Les tests backend sont écrits avec Pytest. Ils utilisent une base de test afin de ne pas modifier les
données de développement. Les fichiers identifiés couvrent notamment :

\begin{itemize}
    \item \texttt{test\_analytics\_dashboard.py} pour les indicateurs du tableau de bord ;
    \item \texttt{test\_assistant\_regions.py} pour les restrictions régionales des assistants ;
    \item \texttt{test\_audit\_logs.py} pour la traçabilité des connexions et des actions ;
    \item \texttt{test\_medicine\_upload.py} pour l'import et la recherche de médicaments.
\end{itemize}

Les tests d'audit vérifient qu'une connexion administrateur réussie crée une entrée
\texttt{ADMIN\_LOGIN}, qu'un échec crée une entrée \texttt{ADMIN\_LOGIN\_FAILED}, et qu'une inscription crée
une entrée \texttt{USER\_REGISTERED}. Les tests de restrictions régionales contrôlent qu'un assistant ne
peut agir que dans son périmètre autorisé. Les tests d'import de médicaments couvrent les colonnes
obligatoires, la normalisation des prix, la gestion des doublons par \texttt{code\_pct} et les messages
d'erreur pour les lignes invalides.

\subsection{Tests API}

Les tests API vérifient les routes publiques et protégées. Ils contrôlent les paramètres requis, les
limites de pagination, les erreurs de validation, les accès interdits et la cohérence des réponses JSON.

\subsection{Validation du déploiement local}

La version validée du projet est exécutable localement avec trois services séparés : backend FastAPI,
administration React/Vite et application mobile Expo. Cette structure reflète l'architecture cible et
donne un cadre clair pour vérifier les échanges API avant un déploiement final.

\begin{table}[H]
    \centering
    \caption{Preuves de lancement et de build}
    \begin{tabular}{L{4cm}L{4.8cm}L{4.2cm}}
        \toprule
        \textbf{Composant} & \textbf{Commande de vérification} & \textbf{État} \\
        \midrule
        Backend FastAPI & \texttt{pytest} sur tests ciblés & Tests ciblés réussis. \\
        Administration web & \texttt{npm run build} & Build de production réussi. \\
        Application mobile & \texttt{npm test -- --runInBand} & Tests Jest réussis. \\
        Configuration & Variables d'environnement listées & Paramètres identifiés. \\
        \bottomrule
    \end{tabular}
\end{table}

\subsection{Validation fonctionnelle}

Les principaux scénarios fonctionnels vérifiés sont présentés dans le tableau suivant.

\begin{longtable}{L{4cm}L{5cm}L{5cm}}
    \caption{Scénarios fonctionnels testés}\\
    \toprule
    \textbf{Scénario} & \textbf{Étapes} & \textbf{Résultat attendu} \\
    \midrule
    \endfirsthead
    \toprule
    \textbf{Scénario} & \textbf{Étapes} & \textbf{Résultat attendu} \\
    \midrule
    \endhead
    Connexion administrateur & Saisir email et mot de passe, puis valider. & Accès au tableau de bord et stockage des jetons. \\
    Recherche mobile & Saisir un nom, une adresse ou un gouvernorat. & Affichage d'une liste de pharmacies pertinentes. \\
    Recherche par proximité & Autoriser la localisation et lancer la recherche. & Pharmacies triées par distance. \\
    Consultation des gardes & Choisir une date dans le calendrier. & Affichage des pharmacies de garde correspondantes. \\
    Import pharmacies & Charger un fichier CSV valide. & Données importées et bilan affiché. \\
    Import gardes & Charger un planning CSV. & Gardes consultables côté administration et côté mobile. \\
    Recherche médicament & Chercher par nom, DCI ou code PCT. & Résultats pertinents et détail accessible. \\
    Audit logs & Filtrer les actions enregistrées. & Résultats paginés cohérents. \\
    Changement de langue & Passer du français à l'arabe. & Interface traduite avec direction RTL appliquée. \\
    \bottomrule
\end{longtable}

\section{Difficultés rencontrées et solutions adoptées}

Plusieurs difficultés sont apparues pendant la réalisation. Elles ont principalement concerné la cohérence
des données, les droits d'accès et l'intégration entre les trois parties du projet.

\begin{longtable}{L{5cm}L{9cm}}
    \caption{Difficultés rencontrées et solutions adoptées}\\
    \toprule
    \textbf{Difficulté} & \textbf{Solution adoptée} \\
    \midrule
    \endfirsthead
    \toprule
    \textbf{Difficulté} & \textbf{Solution adoptée} \\
    \midrule
    \endhead
    Alignement entre backend, web et mobile & Centralisation des contrats API et séparation des services backend. \\
    Qualité variable des fichiers CSV & Validation des colonnes, contrôle des types, retour des erreurs ligne par ligne et comportement d'upsert pour les médicaments. \\
    Gestion des rôles et des assistants régionaux & Vérification des permissions côté backend et application du périmètre régional dans les services. \\
    Géolocalisation mobile & Demande explicite des permissions, contrôle de l'activation des services de position et recherche par coordonnées. \\
    Maintenabilité du backend & Séparation entre routes, services, modèles, schémas, sécurité et migrations. \\
    Internationalisation mobile & Utilisation de fichiers de traduction et prise en charge de la direction RTL pour l'arabe. \\
    \bottomrule
\end{longtable}

\section{Conclusion}

La réalisation a abouti à une plateforme fonctionnelle composée d'une API sécurisée, d'un tableau de bord
administrateur et d'une application mobile. Les tests réalisés ont couvert les parcours essentiels :
authentification, recherche de pharmacies, consultation des gardes, imports CSV, gestion des rôles, audit
et recherche de médicaments. Des améliorations restent possibles, notamment les tests bout en bout, les
tests de charge et la validation sur un volume de données plus important.
